% !TEX program = xelatex

\documentclass{resume}
%\usepackage{zh_CN-Adobefonts_external} % Simplified Chinese Support using external fonts (./fonts/zh_CN-Adobe/)
%\usepackage{zh_CN-Adobefonts_internal} % Simplified Chinese Support using system fonts

\begin{document}
\pagenumbering{gobble} % suppress displaying page number

\name{Zhiyu Ding}

\basicInfo{
  \email{dingzhiyu2004@163.com} \textperiodcentered\ 
  \phone{(+86) 18765788600} \textperiodcentered\ 
  \homepage[nevergpdzy.com]{https://nevergpdzy.com}}

\section{Education}
\datedsubsection{\textbf{Southwest Petroleum University}, Chengdu, Sichuan}{September 2023 -- Present}
\textit{Undergraduate Student} in Data Science and Big Data Technology, School of Computer and Software Engineering
\begin{itemize}
  \item \textbf{Academic Performance}: GPA: 4.13/5.0, Major Ranking: 1/66
  \item \textbf{Languages}: English CET-6 (478)
  \item \textbf{Core Courses}: Advanced Mathematics II(95), Linear Algebra(91), Probability and Statistics(93), Principles of Statistics(93), Data Structures and Algorithms(90)
\end{itemize}

\section{Project Experience}

\subsection{\textbf{Parallel Computing Optimization for Oil Spill Prediction Model}}
\textit{May 2024 -- August 2024}
\role{2024 Marine Computing Challenge Finals}{Team Leader}
Selected for the 2024 Marine Computing Challenge finals; led hybrid parallelization and coastline collision detection optimization for a 2D oil spill prediction model.
\begin{itemize}
  \item Completed MPI+OpenMP hybrid parallelization on the original serial program, analyzing bottlenecks with VTune and optimizing via load balancing, memory access reordering, non-blocking communication, packed communication, segment rapid elimination, and binary search
  \item Passed the organizer's correctness verification, achieved a 2482.14x speedup over the baseline, ranked 5th in the finals, and won the national third prize
\end{itemize}

\subsection{\textbf{Tecorigin Deep Learning Operator Performance Optimization}}
\textit{June 2024 -- December 2024}
\role{2nd OpenAtom Competition -- Tecorigin Operator Development Challenge}{Team Leader}
Selected for the 2nd OpenAtom Competition Tecorigin Operator Development Challenge; conducted systematic performance optimization of Tecorigin convolution forward operators.
\begin{itemize}
  \item Built performance profiles using Perf Data, identifying I/O as the core bottleneck at 93.1\%; designed a 234KB SPM management strategy and double-buffering asynchronous pipeline, and implemented SIMD data reordering based on floatv16 vector registers
  \item Reduced total execution time from 1820.78ms to 489.18ms, achieving a 3.7x overall speedup; resolved the mismatch between matrix multiplication library output and NHWC format, winning the national third prize
\end{itemize}

\subsection{\textbf{PCG Algorithm Optimization on New Generation Sunway Supercomputer}}
\textit{February 2024 -- April 2024}
\role{7th Domestic CPU Parallel Application Challenge}{Individual Project}
Selected for the 7th Domestic CPU Parallel Application Challenge preliminary round; optimized the Preconditioned Conjugate Gradient (PCG) solver for the Sunway heterogeneous many-core architecture.
\begin{itemize}
  \item Restructured the PCG workflow based on the Sunway 64-core slave-core architecture, parallelizing core operators such as SpMV, dot product, and preconditioning using athread, and optimizing data access and blocking strategies via DMA and LDM
  \item Through over 40 iterations, reduced runtime from 1287 seconds to 32.5 seconds, achieving a 39.6x speedup
\end{itemize}

\section{Research Experience}
\subsection{\textbf{Image Processing and Parallel Computing Laboratory (IPPC Lab)}, Southwest Petroleum University}
\textit{May 2024 -- September 2025}
\role{3D Ultrasound Speckle Tracking Algorithm Optimization}{Sichuan Student Innovation Program; Advisor: Prof. Bo Peng}
\begin{itemize}
  \item Rebuilt the serial MATLAB workflow as a CUDA C++ parallel pipeline for high-volume cross-correlation in 3D ultrasound speckle tracking, decomposing it into RF integral volume construction, cross-correlation integral volume generation, and 3D subpixel displacement estimation GPU stages
  \item Optimized memory access for large-scale 3D prefix-sum computation and introduced warp-level cooperation, accelerating the core kernel by 11.11x and reducing end-to-end runtime from 16.32s to 0.22s, about 73x faster
\end{itemize}

\subsection{\textbf{CUDA Acceleration for Large-Scale Image Pretraining Data Generation}}
\textit{October 2025 -- March 2026}
\begin{itemize}
  \item Rewrote large-image tiling, noise injection, and training sequence generation for Hierarchical Denoising Encoder (HDE) pretraining in CUDA C++, packaging the workflow as a Python-callable module
  \item Decomposed the pipeline into integral image computation, quadtree partitioning, hierarchical noise generation, and patch synthesis, processing 8192$\times$8192 images on A100 in 3.16s with 13.0x speedup over CPU, 2.35x over CuPy, and 69.1x faster noise generation
\end{itemize}

\section{Awards and Honors}
\datedline{\textit{International Second Prize}, ASC2025 World Student Supercomputer Challenge}{June 2025}
\datedline{\textit{Invited Participant}, SC24 International Supercomputing Competition Online Track IndySCC}{November 2024}
\datedline{\textit{National Third Prize}, 2024 Marine Computing Challenge Finals}{August 2024}
\datedline{\textit{National Third Prize}, Tecorigin Operator Development Challenge Finals}{December 2024}
\datedline{\textit{Provincial Second Place}, Tianyi Cloud Xirang Cup College AI Competition, Sichuan Provincial Competition}{July 2025}
\datedline{\textit{National Third Prize}, 15th Blue Bridge Cup National Finals}{June 2024}
\datedline{\textit{First-class Scholarship, Second-class Scholarship}, Southwest Petroleum University Outstanding Student Scholarship}{2024}

\section{Additional Information}
\begin{itemize}[parsep=0.5ex]
  \item Blog: https://nevergpdzy.com
  \item GitHub: https://github.com/NeverGpDzy
  \item Research Interests: High Performance Computing and Parallel Programming, GPU/CPU heterogeneous computing optimization
\end{itemize}

%% Reference
%\newpage
%\bibliographystyle{IEEETran}
%\bibliography{mycite}
\end{document}
